% !TeX root = ../main.tex

\chapter{Discussion}\label{chapter:discussion}

No multi-agent configuration improves on the strongest single-agent baseline by more
than $0.011$ pooled $F_2$ points across eleven models.

\section{The Headline Finding and Its Conditions}\label{sec:disc:headline}

On the 102-contract CUAD test split under the upsampled-positive class
distribution described in \S\ref{sec:meth:sampling}, no multi-agent configuration
improves on the strongest single-agent baseline by more than $0.011$ points of
mean $F_2$ pooled across the eleven-model panel.
Under the pre-registered decision rule of Section~\ref{sec:meth:hypotheses},
H1 is rejected in its universal form: no multi-agent configuration achieves
consistent superiority across the eleven-model panel, and eleven of the 22
per-model contrasts are negative or zero. Conditional improvement exists for
individual (model, configuration) cells --- \cfgCX\ (M0) on the OpenAI family
and lighter Anthropic models --- but the pooled signal is dominated by the
model-family interaction. H2 is rejected: both pipeline decomposition
(\cfgCX, M0) and specialist routing (\cfgMA, M1) widen the rare--common $F_2$
gap relative to \cfgZS\ (B0). H3 is supported per model --- \cfgMA\ (M1) and
\cfgCP\ (M2) are statistically distinguishable on every model after BH
correction --- but the direction is inconsistent across the panel and the
pooled $F_2$ difference is zero to three decimals.

The single most informative pattern is the model-family interaction.
The Anthropic family achieves its peak under \cfgCoT\
(B1, $\overline{F_2}{=}0.851$); the Google family under \cfgZS\
(B0, $0.826$); the OpenAI family under \cfgCX\
(M0, $0.783$). The capability-tier gradient in
Figure~\ref{fig:tier-heatmap} is consistent with this: light- and mid-tier
models gain from M0's calibration, while heavy-tier models are better served
by B0, with the $F_2$ capability gap between Heavy and Light tiers narrowing
from $0.106$ under B0 to $0.054$ under M0 --- a $49\%$ reduction. Decomposition acts partly as a capability-equalising effect, but at a cost
(Section~\ref{sec:disc:cascade}) that is not recovered for heavy models
that already absorb long-context load in a single call. Prior work has observed
that a well-prompted single agent matches multi-agent decomposition under
stronger base models \cite{gao2025single}; the present data reproduce this:
the Anthropic Opus tier achieves its best $F_2$ under B1, not under any
decomposed architecture.

The implication for the open empirical question of when multi-agent
decomposition is justified is that, on this benchmark, the conditions are
narrower than the literature on agentic legal AI implies. A practitioner
asking whether decomposition helps a CUAD-class extraction task must answer
two prior questions: which model family, and which capability tier. On the
eleven models tested, only the OpenAI family receives a meaningful
multi-agent benefit at the pooled level, and even that benefit is dominated
in $F_2$/\$ by a B0 deployment on a cheaper Google Flash model.

\section{The Cascade as a Recall--Precision Exchange}\label{sec:disc:cascade}

Section~\ref{sec:res:cascade} identifies the classification-gate cascade as
the dominant failure mode of \cfgCX\ (M0) on the Google family:
$97.8$--$97.9\%$ of M0 false negatives on Gemini~2.5~Flash and Pro originate
at the gate rather than at the extractor. The discussion here defends the
framing of this as a recall--precision exchange rather than as pure regression.

The classifier reduces over-extraction substantially: mean $\omega$ drops
from $0.320$ under \cfgZS\ (B0) to $0.130$ under \cfgCX\ (M0), a $59\%$
reduction (Table~\ref{tab:omega}). Under $F_1$, which weights precision and
recall equally, the gate's gain is sufficient to push M0 to the highest pooled
score ($\overline{F_1}{=}0.798$ vs.\ B0's $0.737$). Under $F_2$, the gain
is insufficient to compensate the recall loss. The cascade is thus not evidence that M0 is broken, but evidence that M0 makes
a different operating-point trade. Under the pre-registered primary metric
$F_2$, the trade is unfavourable for 7 of 11 models. The practitioner question
is which operating point matches the deployment cost structure.

The cascade direction is predictable from a single-agent calibration pass.
Models with high B0 over-extraction rates (Gemini~2.5~Pro
$\omega_{\text{B0}}{=}0.456$; Gemini~2.5~Flash $0.392$) inherit a permissive
prior on negative instances; the gate amplifies the cautious prior,
suppressing recall. Models with low B0 over-extraction (GPT-4.1~Nano
$\omega_{\text{B0}}{=}0.085$) gain from the gate because their single-agent
extraction loses ground to laziness rather than to over-extraction. The
practitioner heuristic is: run B0, inspect $\omega_{\text{B0}}$, and deploy
M0 only if $\omega_{\text{B0}}$ is moderate-to-low. Within the present
eleven-model panel, models with $\omega_{\text{B0}} \lesssim 0.20$ gain on
$F_1$ under M0; those above that threshold do not. The boundary is indicative,
not validated.

This result connects to the failure-mode taxonomy of \textcite{cemri2025why},
who identify the inter-agent verification gap as a leading cause of
multi-agent regression: the downstream extractor cannot override the upstream
classifier's silent rejection. A reflective verifier agent --- absent in
M0 --- would be the natural mitigation, and is identified as an untested
architectural pattern in Section~\ref{sec:concl:outlook}.

\section{Routing as a Deterministic Lookup}\label{sec:disc:routing}

The \cfgMA\ (M1) router agreed with a static category-to-specialist mapping on
$99.91\%$ of requests across the panel (Section~\ref{sec:res:h3}).
The natural reading is that routing in this design space reduces to a table
lookup keyed on the category name, a negative finding for the LLM-router
subspace illustrated by ChatLaw~v2~\cite{cui2024chatlaw} and
L-MARS~\cite{wang2025lmars}. Two qualifications are worth stating.

First, the result depends on three architectural commitments that do not hold
in every multi-agent design: the specialist domains were disjoint and
descriptively named; the categories carried indicators the router could read
directly from the category label; and the partition was fixed before any
model was run. In designs where domains overlap, where the category
vocabulary is open-ended, or where the routing decision must integrate
retrieved context, the router's work becomes non-trivial and routing accuracy
falls below ceiling. The generalisable claim is therefore narrower: in a
closed-vocabulary specialist-routing design, the router's contribution is
replicable by a deterministic dictionary at zero LLM cost.

Second, \cfgMA\ (M1) and \cfgCP\ (M2) differ statistically on every model
despite identical routing. The difference is not in dispatch but in prompt
format: M1 sees one focused specialist preamble per call, M2 sees the textual
union of three.
Consolidating the specialist content into a single prompt produces measurably
better span coverage (Jaccard $+0.071$ pooled, Table~\ref{tab:h3}) without
changing $F_2$. The mechanism is plausibly attentional: the union prompt
presents all three category-cluster framings simultaneously and lets the
model choose, rather than committing the dispatch decision before the
contract is read.

The combined implication for the design of legal-AI multi-agent systems is
that the value of specialist decomposition lies in the prompt content, not
in the dispatch architecture. A practitioner with a fixed specialist
taxonomy is better served by a single agent receiving the union of
specialist instructions than by a router-plus-specialists topology, and
eliminates the $23\%$ orchestrator cost premium.

\section{The $\beta{=}2$ Choice and Operating-Point Sensitivity}\label{sec:disc:metric}

The pooled architectural ranking depends on the chosen $\beta$. At
$\beta{=}1$ the ranking is \cfgCX\ $>$ \cfgMA\ $>$ \cfgZS\ $>$ \cfgCP\ $>$ \cfgCoT.
At $\beta{=}2$ the ranking is \cfgZS\ $>$ \cfgCoT\ $>$ \cfgCX\ $>$ \cfgMA\ $\approx$ \cfgCP. The architectural
conclusion of the thesis therefore depends, in part, on the $\beta$ choice.

The commitment to $\beta{=}2$ is justified in
Section~\ref{sec:meth:metrics} by the deployment cost structure of CLM
review: a missed clause carries the legal cost of an undisclosed obligation,
whereas a spurious extraction is corrected at the lower cost of attorney
review time. The empirical question of where the cost ratio actually sits is
deployment-dependent and falls outside the scope of this thesis. The
sensitivity is mitigated by reporting $F_1$, recall, precision, $\omega$,
and FN-rate alongside $F_2$, so a deployment with a different cost ratio can
re-rank architectures from the same data. The practitioner recommendation is
conditional on $\beta$: at $\beta{=}2$ deploy B0; at $\beta{=}1$ deploy M0;
below $\beta{=}1$ (where false positives are costlier than false negatives,
as in some compliance-screening workflows), M0's calibration advantage
widens further.

\section{Benchmark-to-Production Gaps}\label{sec:disc:practitioner}

CUAD measures one axis --- extraction accuracy under a fixed forty-one-category
taxonomy on clean born-digital text with expert ground truth --- that is
necessary but not sufficient for production deployment. Three deployment axes
relocate the architectural argument.

\paragraph{Preprocessing and retrieval are upstream of architecture.}
CUAD contracts are EDGAR filings: clean, machine-readable, consistently
formatted. Enterprise repositories contain scanned PDFs of variable OCR
quality, redactions, and partial digitisations; document-preprocessing
quality is the primary performance ceiling, set before any extraction
architecture is invoked. CUAD's one-to-one (contract, category) framing
also omits the retrieve-then-extract pipeline production reviewers actually
run --- which contracts contain uncapped liability across the repository?
The classification-gate cascade (\S\ref{sec:disc:cascade}) is the
single-contract analogue of that retrieval failure: the same
upstream-rejection dynamics that suppress M0 recall on Gemini~2.5 will
plausibly suppress retrieval-extraction recall in production.

\paragraph{Open taxonomies and cross-contract reasoning are unmodelled.}
CUAD's forty-one categories cover general commercial provisions; construction
contracts (variation orders, retention, defects liability), SaaS agreements
(SLA penalties, sub-processor notice), and jurisdictional clause variants
do not map cleanly onto the schema. The fixed three-domain specialist
partition tested here (M1, M2) is not extensible to an open taxonomy without
re-routing. Cross-contract questions --- consistency of liability caps
across vendor agreements; impact analysis on a counterparty acquisition ---
require portfolio-level orchestration over per-contract agents that no
public benchmark currently evaluates and that this thesis identifies as
future work (\S\ref{sec:concl:outlook}).

\paragraph{Production has no ground truth, so trace primitives carry the value.}
A production deployment cannot compute $F_2$, Jaccard, or $\omega$;
evaluation proceeds through user feedback, spot-checking, and exception
triage. The most valuable architectural property in that regime is whether
each decision is attributable to a specific agent, prompt, and routing
branch (\S\ref{sec:auditability}). The M0 classifier verdict and the M1
routed-specialist label are exactly the trace primitives a reviewer needs,
together with the schema-enforced $1.000$ grounding rate of M1
(Table~\ref{tab:grounding}). The benchmark $F_2$ ranking therefore does not
translate into ``deploy single-agent in production'': the production utility
function includes auditability, modular update boundaries, and integration
seams that the benchmark does not measure. This is a defensible reason to
build multi-agent CLM systems --- weaker than ``decomposition wins on
accuracy'', but cleanly stated.

\section{Threats to Validity}\label{sec:disc:threats}

The threats formalised in Section~\ref{sec:meth:threats} bound the
generalisation of these findings. Three are most consequential.

\paragraph{Single benchmark.}
The evaluation is on the 102-contract CUAD test split under the upsampled-positive
class distribution (\S\ref{sec:meth:sampling}). The findings may not extrapolate
to non-CUAD clause taxonomies, to non-English contracts, or to non-U.S.\
jurisdictions.
The contract-length moderation result
(Section~\ref{sec:res:length}, Table~\ref{tab:length-quartile}) shows that
the architectural ranking is not stable across document lengths within CUAD
itself: multi-agent helps on short contracts and hurts on long ones.
The architectural ranking observed here is therefore a ranking on this
evaluation envelope, not a universal ranking.

\paragraph{Pinned model snapshots.}
Model identifiers are pinned to specific snapshots in the run metadata.
Provider behaviour can shift under silent model updates; the family
interaction observed here may not be stable across future provider releases.
The qualitative pattern --- that architecture choice interacts with model
family more than with any within-architecture variable --- is consistent with
findings reported across task domains in the recent literature
\cite{gao2025single,cemri2025why}, which suggests the pattern is consistent across task domains even if specific
point estimates are not.

\paragraph{Closed proprietary panel.}
Open-weights models are excluded (Section~\ref{sec:meth:models}).
The proprietary--open-source gap reported by \textcite{liu2025contracteval} is
not re-examined here. The finding that multi-agent decomposition does not
improve pooled $F_2$ may not transfer to a panel of weaker open-weights
models, where decomposition's capability-equalisation effect
(Section~\ref{sec:disc:headline}) might be larger. Open-weights extension
is identified in Section~\ref{sec:concl:outlook} as the most consequential
extension of the present work.
